\documentstyle[% 


righttag% 


,11pt% 


,amssymb% 


,russian%               remove in Eng. 


,izvru%                 Set izven% in Eng. 


]{amsart} 


 


% {,}


\newcommand{\tts}[1]{} 


 


\theoremstyle{definition} 


\theorembodyfont{\rm} 


\newtheorem{notenonum}{\hskip\parindent Замечание} 


\renewcommand{\thenotenonum}{} 


\numberwithin{equation}{section} 


 


%\hoffset=-15mm%                  For Eng. 


\setcounter{page}{73} 


\god{2004} 


\nomer{9~(508)} %                        Remove in Eng. 


%\nomer{Vol.\,48, No.\,9}%               Set in Eng. 


%\str{\pageref{begin}--\pageref{end}}%  Set in Eng. 


\udk{517.955:532.5} 


 


\author{\it Л.Д.\,Эскин} 


% NB:   ^^ remove in Eng. 


\title{К задаче П.Я.\,Полубариновой-Кочиной\\ об опорожнении бассейна} 


 


\date{} 


%\pagestyle{myheadings}%         Set in Eng. 


\pagestyle{plain} %              Remove in Eng. 


\begin{document} 


%\thispagestyle{plain}%          Set in Eng. 


\maketitle 


\label{begin} 


 


 


\setcounter{section}{1} 


 


 


{\bf 1.} Уравнение нелинейной теплопроводности 


\begin{equation} 


u_t=(u^m u_x)_x 


\end{equation} 


возникает при моделировании самых разнообразных процессов в механике 


сплошной среды. Сюда относятся процессы переноса тепла (массы) в среде с 


коэффициентом теплоемкости (диффузии), степенным образом зависящим от 


температуры (концентрации), фильтрации политропного газа в пористую среду. 


В последнем случае уравнение (1.1) называется уравнением Лейбензона, а при 


$m=1$ --- уравнением Буссинеска. Это же уравнение возникает и при 


моделировании динамики грунтовых вод. Уравнение (1.1) будет 


рассматриваться на полуоси $x>0$ вместе с начальным и граничным условиями 


\begin{equation} 


u(x,0)=1,\quad u(0,t)=0 


\end{equation} 


(известная задача П.Я.\,Полубариновой-Кочиной [1] об опорожнении 


бассейна). Решение задачи (1.1), (1.2) при $m=1$ было протабулировано в 


[1], а еe приближенное аналитическое решение построено в [2]. Цель данной 


статьи --- построение асимптотики решения задачи (1.1), (1.2) вблизи 


границ $x=0$ и $x=\infty$. Поскольку предлагаемая здесь методика связана 


с качественной теорией динамических систем и теорией нормальных форм 


обыкновенных дифференциальных уравнений, будем предполагать, что $m$ --- 


натуральное число. Задача (1.1), (1.2) рассматривается в п.\,2--5. В п.\,6 


аналогичная методика используется для изучения процесса турбулентной 


фильтрации политропного газа в пористой среде. 


\medskip 


\addtocounter{section}{1} 


\setcounter{equation}{0}


 


{\bf 2.} Известно [1], [2], что решение задачи (1.1), (1.2) автомодельно и 


монотонно возрастает по $x$. Полагая $\xi=xt^{-1/2}$, $u=f(\xi)$, из 


(1.1), (1.2) получим для $f(\xi)$ обыкновенное дифференциальное уравнение 


\begin{equation} 


(f^m f_\xi)_\xi+2^{-1}\xi f_\xi=0 


\end{equation} 


и граничные условия 


\begin{equation} 


f(0)=0,\quad f(\infty)=1. 


\end{equation} 


Следуя [3], [4], положим 


\begin{equation} 


f=\xi^{2/m}\phi(\xi),\quad \sigma=\xi \phi_\xi. 


\end{equation} 


В результате подстановки (2.3) в (2.1) будем иметь обыкновенное 


дифференциальное уравнение для $\sigma=\sigma(\phi)$: 


\begin{gather} 


\sigma_\phi=Q/P,\quad Q=((1+2/m)\phi^m+m\phi^{m-1} 


\sigma+2^{-1})w+2m^{-1}\phi^m\sigma,\\ 


P=-\phi^m\sigma,\quad w=\sigma+2m^{-1}\phi,\notag 


\end{gather} 


причем $f_\xi=\xi^{-1+2/m}w$. Уравнение (2.4) удобно заменить динамической 


системой 


\begin{equation} 


\phi_\tau=P,\quad \sigma_\tau=Q. 


\end{equation} 


Выбор независимой переменной $\tau$ роли не играет, т.\,к. нас будут 


интересовать лишь траектории динамической системы (2.5) в фазовой 


плоскости $(\phi,\sigma)$, а не закон движения по траектории. Отметим 


прежде всего, что, имея интегральную кривую уравнения (2.4), можно найти 


из второго уравнения (2.3) однопараметрическое семейство его решений 


$\phi=\phi(\xi,c)$, а затем с помощью первого из равенств (2.3) --- 


однопараметрическое семейство автомодельных решений уравнения (1.1). 


Следствием неотрицательности и монотонного возрастания по $x$ решения $u$ 


задачи (1.1), (1.2) являются неотрицательность и монотонное возрастание 


по $\xi$ функции $f(\xi)$. Поэтому будем рассматривать лишь траектории 


динамической системы (2.5), принадлежащие сектору $I$ фазовой плоскости 


$(\phi,\sigma)$, который лежит в полуплоскости $\phi>0$ и ограничен 


полуосями $\phi=0$ ($\sigma>0$) и $w=0$ ($\phi>0$). С помощью стандартной 


методики качественной теории динамических систем на плоскости [5] нетрудно 


исследовать поведение траекторий уравнения (2.5) в секторе $I$. 


Ограничимся тем, что приведем лишь необходимые для дальнейшего 


окончательные результаты проведенного анализа (см. рис.\,1). 


\medskip 


 


\begin{center} 


\unitlength=1mm 


\begin{picture}(170,86) %%170= width, 73=heigth, added 5 mm for signature 


\put(50,86){\special {em:graph eskin6.bmp}} %% left X, upper Y, -, /2, +toX 


\put(78,0){\mbox Рис.\,1} % added 5mm (73+5) 


\end{picture} 


\end{center} 


\medskip 


 


\noindent{}Сектору $I$ принадлежит ветвь изоклины нуля уравнения (2.4) 


\begin{gather*} 


\sigma_1(\phi)=a(\phi)+b(\phi),\quad a(\phi)=-((3+4/m) 


\phi^m+1/2)/(2m\phi^{m-1}),\\ 


b(\phi)=\{(1+4/m)^2\phi^{2m}+(4m^{-1}-1)\phi^m+1/4\}^{1/2}. 


\end{gather*} 


Эта ветвь расположена в четвертом квадранте фазовой плоскости 


$(\phi,\sigma)$, входит в начало координат в направлении с угловым 


коэффициентом $k_1=-2/m$ (т.\,е. касается прямой $w=0$), уходит в $\infty$ 


в направлении с угловым коэффициентом $k_2=-1/m$. С координатными осями и 


прямой $w=0$ эта ветвь пересекается лишь в начале координат. Ветвь 


$\sigma_1$ и координатная полуось $\phi$ ($\phi>0$) разбивают сектор $I$ 


на три подсектора $I_1$, $I_2$, $I_3$, знаки ``$+$', ``$-$'' на рис.\,1 


указывают знаки производной вдоль траектории в этих подсекторах. Все 


траектории в секторе $I$ стремятся к положению равновесия в направлении с 


угловым коэффициентом $k_1$, причем имеется два однопараметрических 


семейства $S_1$ и $S_2$ траекторий. Кривые семейства $S_1$ переходят из 


подсектора $I_3$ в подсектор $I_2$, горизонтально пересекая ветвь изоклины 


нуля, затем вертикально пересекают координатную ось $\phi$, входят в 


подсектор $I_1$ и приближаются на $\infty$ к полуоси $\sigma>0$ 


по асимптотическому закону 


\begin{equation} 


\sigma\sim A\phi^{-m},\quad \phi\to0, 


\end{equation} 


$A>0$ --- параметр семейства. Кривые семейства $S_2$ не пересекают ни 


изоклины нуля, ни прямой $w=0$ и на всем протяжении принадлежат подсектору 


$I_3$. Для кривых семейства $S_2$ справедлива асимптотика при 


$\phi\to\infty$


\begin{equation} \sigma\sim 


-2m^{-1}\phi+B\phi^{\alpha},\quad \alpha=1-2^{-1}m,


\end{equation}


$B>0$ --- параметр семейства. Семейства $S_1$ и $S_2$ разделяет кривая $S$ --- 


ветвь сепаратрисы седла на $\infty$ динамической системы (2.5) в 


направлении с угловым коэффициентом $k_3=\lambda_1=-\beta/m>-2/m$, 


$\beta=(m+2)/(m+1)$. Этот результат нетрудно получить, выполнив 


преобразование Пуанкаре 


\begin{equation} 


\phi=1/y_1,\quad \sigma=(\lambda_1+y_2)/y_1, 


\end{equation} 


которое преобразует систему (2.5) в систему 


\begin{equation} 


y_{1\tau}=\lambda_1y_1+y_1y_2,\ \ y_{2\tau}=\lambda_2y_2+ 


(m+1)(y_2)^2+(2(m+1))^{-1}(y_1)^m+2^{-1} 


(y_1)^my_2, 


\end{equation} 


$\lambda_2=1$, причем седло на $\infty$ системы (2.5) преобразуется в 


седло $y_1=y_2=0$ ($\lambda_1\lambda_2<0$) системы (2.9), а кривая $S$ --- 


в одну из ветвей сепаратрисы этого седла. С помощью асимптотических формул 


$$ 


\sigma\sim -2m^{-1}\phi\ \ (\phi\to0),\quad 


\sigma\sim \lambda_1\phi\ \ (\phi\to\infty) 


$$ 


для кривой $S$ и аналогичных асимптотических формул (2.6), (2.7) для 


кривых семейств $S_1$ и $S_2$ с учетом соотношений (2.3) нетрудно показать,


что только однопараметрическое семейство решений $f(\xi)$ уравнения 


(2.1), порожденное в силу (2.3) кривой $S$, удовлетворяeт граничному 


условию (2.2) при $\xi=0$ и условию $\lim f(\xi)=c<\infty$ при 


$\xi\to\infty$. При этом асимптотическое разложение $f(\xi)$ при $\xi\to0$ 


определяется асимптотикой кривой $S$ при $\phi\to\infty$, т.\,е. 


асимптотикой в окрестности седла на $\infty$, а асимптотическое разложение 


$f(\xi)$ при $\xi\to\infty$ --- асимптотикой $S$ при $\phi\to0$, т.\,е. 


асимптотикой в окрестности положения равновесия $\phi=0$, $\sigma=0$. 


После получения этих асимптотик надо будет выбрать значение параметра 


семейства решений, порожденных в силу (2.3) кривой $S$, таким образом, чтобы 


выполнялось и граничное условие (2.2) на $\infty$. 


\medskip 


\addtocounter{section}{1} 


\setcounter{equation}{0}


 


{\bf 3.} Для вычисления асимптотического разложения для кривой $S$ при 


$\phi\to\infty$ (в окрестности седла на $\infty$) приведем систему (2.9) к 


нормальной форме в окрестности седла $y_1=y_2=0$. Система (2.9) при 


$m\ge2$ имеет канонический вид, и в этом пункте будем предполагать это 


условие выполненным. Случай $m=1$ рассмотрим в п.\,4. Будем использовать 


обозначения и результаты из [6], полагая $Q=(d,e)$ $Z^Q=z^d_1z^e_2$. В 


([6], с.\,98) показано, что если $\lambda=\lambda_1/\lambda_2=-r/s$, где 


$r/s$ --- несократимая рациональная дробь, то система (2.9) с помощью 


преобразования 


\begin{equation} 


y_i=z_i+h_i(Z),\quad i=1,2,


\end{equation} 


где $h_i(Z)=\sum\limits_Q h_{iQ}Z^Q$ --- степенной ряд, не содержащий 


свободного и линейных 


слагаемых, в окрестности седла $y_1=y_2=0$ приводится к нормальной форме 


\begin{equation} 


z_{i\tau}=z_i\bigg(\lambda_i+\sum^\infty_{k=1} 


g_{i(ks,kr)} z^{ks}_1z^{kr}_2\bigg). 


\end{equation} 


В случае системы (2.9) и нечетного $m$ имеем $r=m+2$, $s=m(m+1)$, а в 


случае четного $m=2t$ имеем $r=t+1$, $s=t(2t+1)$. Седло $y_1=y_2=0$ 


преобразование (3.1) переводит в седло $z_1=z_2=0$ системы (3.2). 


Очевидно, сепаратриса этого седла имеет ветвь $z_1=\exp(\lambda_1\tau)$, 


$z_2=0$ (она входит в седло при $\tau\to\infty$). С учетом преобразования 


(3.1) и того обстоятельства, что степенные ряды $h_i(Z)$ не содержат 


свободного и линейных слагаемых, получаем для седла $y_1=y_2=0$ ветвь 


сепаратрисы 


\begin{equation} 


y_1=\exp(\lambda_1\tau)+\sum^\infty_{k=2}h_{1(k,0)} 


\exp(\lambda_1k\tau),\quad 


y_2=\sum^\infty_{k=2}h_{2(k,0)}\exp(\lambda_1k\tau). 


\end{equation} 


Коэффициенты $h_{i(k,0)}$ могут быть последовательно определены в 


результате подстановки формальных рядов (3.3) в систему (2.9) и применения 


метода неопределенных коэффициентов. Оказывается, что $y_1$, $y_2$ 


разлагаются в ряды по степеням $u=\exp m\lambda_1\tau$. А именно, 


\begin{equation} 


y_1=u^{1/m}\bigg(1+\sum^\infty_{k=1}a_k u^k \bigg),\quad 


y_2=\sum^\infty_{k=1}b_k u^k, 


\end{equation} 


где 


\begin{gather} 


a_1=(2\beta(2m+3))^{-1},\quad b_1=-(2(2m+3))^{-1},\\ 


a_2=(m+1)^2/(8\beta^2(2m+3)^2(3m+5)),\notag\\ 


b_2=-(m^2-m-4)/(4\beta(2m+3)^2(3m+5)),\notag 


\end{gather} 


а остальные коэффициенты $a_k$, $b_k$ определяются из рекуррентных 


соотношений, которые ввиду их громоздкости не приводятся. Соотношения 


(3.4) задают параметрическое представление ветви сепаратрисы седла 


$y_1=y_2=0$ системы (2.9), причем малой окрестности седла соответствуют 


малые значения параметра $u$. Полагая 


$$ 


\bigg(1+\sum^\infty_{k=1}a_k u^k \bigg)^{-1}=1+ 


\sum^\infty_{k=1}r_k u^k, 


$$ 


где коэффициенты $r_k$ выражаются через $a_s$ ($s=1,\dotsc,k$) и могут 


последовательно определяться с помощью метода неопределенных 


коэффициентов, с учетом соотношений (2.8) получим параметрическое 


представление ветви сепаратрисы седла на $\infty$ системы (2.5) 


\begin{gather} 


\phi(u)=y^{-1}_1=u^{-1/m}\bigg(1+\sum^\infty_{k=1} 


r_k u^k \bigg),\\ 


\sigma(u)=(\lambda_1+y_2)/y_1=u^{-1/m}\bigg\{\lambda_1+ 


(\lambda_1r_1+b_1)u+\sum^\infty_{k=2}\bigg( 


\lambda_1r_k+b_k+\sum_{s+t=k}r_s b_t\bigg)u^k\bigg\}.\notag 


\end{gather} 


В (3.6) \ $r_1=-a_1$, $r_2=a^2_1-a_2$, для остальных коэффициентов 


получаются все более громоздкие формулы. Из (3.6) без труда находим, что 


при $u\to0$ эта кривая стремится к седлу на $\infty$ системы (2.5) в 


направлении с угловым коэффициентом $k_3=\lambda_1$, т.\,е. эта кривая 


совпадает с кривой $S$ и формулы (3.6) дают искомое асимптотическое 


представление этой кривой в окрестности седла. Обозначим 


\begin{gather*} 


d(u)=\lambda_1u^{-1/m}(\sigma(u))^{-1}-1= 


\sum^\infty_{k=1}d_k u^k,\\ 


p(u)=\sum^\infty_{s=1}p_s u^s=d(u)\sum^\infty_{k=1} 


\beta^{-1}r_k(1-km)u^{k-1}. 


\end{gather*} 


Из второго уравнения (2.3) получим 


\begin{equation} 


\xi=c\exp\int\sigma^{-1}d\phi=c u^{1/\beta}\exp\bigg\{ 


\sum^\infty_{k=1}((1+k)^{-1}p_k u+(d_k+r_k(1-km))/ 


(k\beta))u^k\bigg\}. 


\end{equation} 


Из (3.7) следует, что $\xi\to0$ при $u\to0$. Исключая $u$ из (3.6) и 


(3.7), получим разложение $\phi(\xi)$ по дробным степеням $\xi$, т.\,е. 


асимптотическое разложение $\phi(\xi)$ при $\xi\to0$, а затем с помощью 


первого из равенств (2.3) --- и искомое разложение в ряд по дробным 


степеням $\xi$ однопараметрического семейства решений $f(\xi,c)$ уравнения 


(2.1), удовлетворяющего граничному условию (2.2) при $\xi=0$. Это 


разложение имеет вид 


\begin{equation} 


f(\xi)=c^{2/m}(\xi/c)^{1/(m+1)}\sum^\infty_{k=0}l_k 


(\xi/c)^{k\beta}, 


\end{equation} 


$c>0$ --- параметр семейства. Коэффициенты $l_k$ в (3.8) выражаются через 


коэффициенты $a_s$, $b_s$, $r_s$, $d_s$, $p_s$ с $s\le k$, т.\,е. через 


коэффициенты рядов (3.4). Получим разложение (3.8), попутно вычислим и его 


коэффициенты $l_0$, $l_1$, $l_2$, т.\,е. трехчленную асимптотику $f(\xi)$ 


при $\xi\to0$. Будем иметь 


\begin{gather*} 


d_1=a_1-(b_1/\lambda_1),\ \ d_2=a_2+(b_1/\lambda_1)^2- 


(a_1b_1+b_2)/\lambda_1,\ \ p_1=a_1d_1(m-1)/\beta,\\ 


% 


\exp\bigg\{\sum^\infty_{k=1}((1+k)^{-1}p_k u+(d_k+r_k(1-km)) 


/(k\beta))u^k\bigg\}=1+\sum^\infty_{k=1}m_k u^k,\\ 


% 


m_1=(d_1+(m-1)a_1)/\beta,\ \ m_2=2^{-1}(p_1+m^2_1)+ 


(2\beta)^{-1}(d_2+(1-2m)r_2),\\ 


% 


(\xi/c)^\beta=u\bigg(1+\sum^\infty_{k=1}g_k u^k\bigg),\ \ 


g_1=m_1\beta,\ \ g_2=(m_2+2^{-1}m^2_1/(m+1))\beta,\\ 


% 


u=(\xi/c)^\beta\bigg(1+\sum^\infty_{k=1}v_k(\xi/c)^{k\beta} 


\bigg),\ \ v_1=-g_1,\ \ v_2=2g^2_1-g_2, 


\end{gather*} 


откуда и следует разложение (3.8), причем 


\begin{gather*} 


l_0=1,\quad l_1=r_1-m^{-1}v_1,\\ 


l_2=((m+1)v^2_1-2m v_2)(2m^2)^{-1}+m^{-1}(m-1)r_1v_1+r_2. 


\end{gather*} 


Нетрудно вычислить коэффициенты $l_3$, $l_4$, но для них получаются 


слишком громоздкие формулы, и мы их не приводим. Из приведенных формул 


следует, что коэффициенты $l_1$, $l_2$ выражаются только через 


коэффициенты $a_1$, $b_1$, $a_2$ и, следовательно, могут быть явно 


выражены через $m$ с помощью соотношений (3.5). Значение параметра $c$, 


определяющее решение $f(\xi)$ уравнения (2.1), которое удовлетворяет обоим 


граничным условиям (2.2), находится из условия $f(\infty)=1$. Оно зависит 


от $m$ и для каждого конкретного значения $m$ должно определяться в 


численном эксперименте. Наиболее удобно поступить следующим образом. 


Положим $f_1(\xi)=f(\xi,c)|_{c=1}$. С помощью асимптотики (3.8) нетрудно 


вычислить данные Коши $f_1$, $f_{1\xi}$ в точке $\xi_0\ll1$ (асимптотика 


производной получается дифференцированием асимптотики (3.8)) и, решая 


задачу Коши для уравнения (2.1) с этими начальными данными, вычислить 


$f_1(\infty)$. Полагая $f=f(\xi,c_0)$, где $c_0=(f_1(\infty))^{-m/2}$, 


получим решение $f$ уравнения (2.1), удовлетворяющее обоим граничным условиям 


(2.2). 


 


Докажем, что ряд (3.8) сходится в некоторой окрестности точки $\xi=0$. 


Очевидно, для этого достаточно доказать, что в некоторой окрестности точки 


$u=0$ сходятся ряды (см. формулы (3.4)) 


\begin{gather*} 


x_1=u^{-1}y^m_1=\bigg(1+\sum^\infty_{k=1}a_k u^k\bigg)^m= 


\sum^\infty_{k=0}\zeta_k u^k,\quad \zeta_0=1,\\ 


x_2=u^{-1}y_2=\sum^\infty_{k=0}\mu_k u^k,\quad 


\mu_k=b_{k+1},\ \ \mu_0=-(2(2m+3))^{-1},\\ 


u=\exp(m\lambda_1\tau),\quad \lambda_1=-(m+2)/(m(m+1)). 


\end{gather*} 


Из уравнений (2.9) для $y_1$, $y_2$ получаем уравнения для $x_1$, $x_2$ 


\begin{gather*} 


x_{1u}=-m(m+1)(m+2)^{-1}x_1x_2,\\ 


u x_{2u}=-(m+2)^{-1}((2m+3)x_2+2^{-1}x_1+ 


2^{-1}(m+1)u x_1x_2+(m+1)^2u x^2_2). 


\end{gather*} 


Из этих уравнений для коэффициентов $\zeta_k$, $\mu_k$ получаем 


рекуррентные соотношения 


\begin{gather} 


(k+1)\zeta_{k+1}=-m(m+1)(m+2)^{-1}\sum_{p+s=k}\zeta_p\mu_s,\\ 


((k+1)(m+2)+2m+3)\mu_{k+1}=-\bigg(2^{-1}\bigg(\zeta_{k+1}+(m+1) 


\sum_{p+s=k}\zeta_p\mu_s\bigg)+(m+1)^2\sum_{p+s=k} 


\zeta_p\mu_s\bigg).\notag 


\end{gather} 


Из соотношений (3.9), положив $k=0$, находим 


$\zeta_1=m(m+1)(2(m+2)(2m+3))^{-1}$,\linebreak $\mu_1=-(m+1)(m^2-m-4)(4(3m+5)(m+2) 


(2m+3)^2)^{-1}$, откуда следует справедливость оценок 


\begin{equation} 


|\zeta_1|\le 1/4,\quad |\mu_1|\le(32(m+1))^{-1}. 


\end{equation} 


При любом $k$ из (3.9) получаем 


\begin{gather} 


|\zeta_{k+1}|\le m(m+1)((m+2)(k+1))^{-1}\sum_{p+s=k} 


|\zeta_p|\,|\mu_s|,\\ 


% 


|\mu_{k+1}|\le((m+2)(k+1)+2m+3)^{-1}\bigg(2^{-1}\bigg( 


|\zeta_{k+1}|+\notag\\


% 


+(m+1)\sum_{p+s=k}|\zeta_p|\,|\mu_s|\bigg)+(m+1)^2 


\sum_{p+s=k}|\zeta_p|\,|\mu_s|\bigg). \notag


\end{gather} 


Предположение индукции 


\begin{equation} 


|\zeta_p|\le 1/4,\quad |\mu_p|\le(4(m+1))^{-1},\quad 


p=1,2,\dotsc,k, 


\end{equation} 


в силу оценок (3.10) выполняется при $k=1$, а также и для $\mu_0$. С 


помощью 


предположения индукции (3.12)из первого неравенства в (3.11) получаем 


\begin{align*} 


|\zeta_{k+1}|&\le m(m+1)((m+2)(k+1))^{-1}(|\mu_k|+k/ 


(16(m+1)))\le \\ 


&\le m(k+4)(16(m+2)(k+1))^{-1}\le 1/4, 


\end{align*} 


после чего из второго неравенства в (3.11) следует 


$$ 


|\mu_{k+1}|\le(32((m+2)(k+1)+2m+3))^{-1}(3k+10) 


\le(4(m+1))^{-1}. 


$$ 


Таким образом, по индукции доказана справедливость оценок (3.12) при всех 


$p\ge1$, а из этих оценок следует сходимость рядов (3.4) в некоторой 


окрестности точки $u=0$, следовательно, и ряда (3.8) в некоторой 


окрестности точки $\xi=0$. 


\medskip 


\addtocounter{section}{1} 


\setcounter{equation}{0}


 


{\bf 4.} В случае $m=1$ (уравнение Буссинеска) система (2.9) не является 


канонической, но легко приводится к каноническому виду 


\begin{equation} 


v_{1\tau}=\lambda_1v_1-0{,}1v^2_1+v_1v_2,\quad 


v_{2\tau}=v_2-0{,}04v^2_1+0{,}2v_1v_2+2v^2_2,\quad 


\lambda_1=-3/2,


\end{equation} 


линейным преобразованием 


\begin{equation} 


y_1=v_1,\quad y_2=-0{,}1v_1+v_2. 


\end{equation} 


Повторяя рассуждения предыдущего пункта, найдем, что ветвь сепаратрисы 


седла $v_1=v_2=0$ системы (4.1), в которую переходит кривая $S$ в 


результате преобразований (2.8) и (4.2), допускает параметрическое 


представление 


$$ 


v_1=u\bigg(1+\sum^\infty_{k=1}a_k u^k\bigg),\quad 


v_2=\sum^\infty_{k=2}b_k u^k, 


$$ 


$a_1=1/15$, $a_2=1/900$, $b_2=0{,}01$, $b_3=1/1650$. Остальные 


коэффициенты могут быть определены с помощью рекуррентных соотношений 


\begin{gather*} 


k\lambda_1a_k=-0{,}2a_{k-1}-0{,}1\sum_{s+r=k-1}a_s a_r+b_k 


+\sum_{s+r=k}a_s b_r,\\ 


(k\lambda_1-1)b_k=-0{,}08a_{k-2}-0{,}04\sum_{s+r=k-2} 


a_s a_r+0{,}2b_{k-1}+0{,}2\sum_{s+r=k-1}a_s b_r+ 


2\sum_{s+r=k}b_s b_r 


\end{gather*} 


($a_k$ определяется вслед за $b_k$). Опуская выкладки, аналогичные 


проведенным в п.\,3, укажем, что степенное разложение семейства решений 


уравнения (2.1), удовлетворяющего граничному условию $f(0)=0$, получается 


в виде 


$$ 


f(\xi)=c^2(\xi/c)^{1/2}\bigg(1+\sum^\infty_{k=1}l_k 


(\xi/c)^{3k/2}\bigg), 


$$ 


где $l_1=-1/15$, $l_2=1/300$. Параметр $c$ снова определяется из условия 


$f(\infty)=1$ 


с помощью приема, указанного в конце п.\,3 в случае произвольного $m$. 


\medskip 


\addtocounter{section}{1} 


\setcounter{equation}{0}


 


{\bf 5.} Перейдем к изучению асимптотики решения граничной задачи (2.1), 


(2.2) при $\xi\to\infty$. Динамическую систему (2.5) теперь удобно 


переписать в виде системы 


\begin{equation} 


y_{1\tau}=\lambda_1y_1+y_1f_1,\quad y_{2\tau}= 


\lambda_2y_2+y_2f_2, 


\end{equation} 


где $\lambda_1=0$, $\lambda_2=1/2$, $y_1=\phi$, $y_2=w$, 


\begin{equation} 


f_1=2m^{-1}y^m_1-y^{m-1}_1y_2,\quad 


f_2=(2m^{-1}-1)y^m_1+m y^{m-1}_1y_2. 


\end{equation} 


Система (5.1) каноническая. В ([6], с.\,97) показано, что система вида 


(5.1) с $\lambda_1=0$ и $\lambda_2\ne0$ преобразованием (3.1), где снова 


формальные степенные ряды $h_i(Z)$ не имеют свободного и линейных 


слагаемых, приводится к нормальной форме 


\begin{equation} 


z_{1\tau}=z_1g_1(z_1),\quad z_{2\tau}=z_2g_2(z_1), 


\end{equation} 


$g_1(z_1)=\sum\limits^\infty_{k=1}g_{1(k,0)}z^k_1$, $g_2(z_1)=\lambda_2+ 


\sum\limits^\infty_{k=1}g_{2(k,0)}z^k_1$. Обозначим $h_i(Z)=z_i 


h^{(i)}(Z)$. Коэффициенты $g_{i(k,0)}$ степенных рядов (5.2) и 


коэффициенты $h_{iQ}$ ($Q=(q_1,q_2)$) степенных рядов 


$$ 


h^{(1)}(Z)=\sum_Q h_{1Q}Z^Q\ \ (q_1\ge-1,\ \ q_2\ge0),\quad 


h^{(2)}(Z)=\sum_Q h_{2Q}Z^Q\ \ (q_1\ge0,\ \ q_2\ge-1) 


$$ 


последовательно определяются из соотношений ([6], с.\,95) 


\begin{gather} 


h_{i(k,0)}=0,\quad g_{i(k,s)}=0\ \ (s\ne0),\quad 


g_{i(k,0)}=c^{(1)}_{i(k,0)}+c^{(2)}_{i(k,0)},\\ 


h_{iQ}=2q^{-1}_2(c^{(1)}_{iQ}+c^{(2)}_{iQ})\ \ (q_2\ne0),\notag 


\end{gather} 


где 


\begin{equation} 


c^{(1)}_{iQ}=-\sum_{P+R=Q}h_{iP}g_{iR}-\sum_{P+R=Q} 


(p_1g_{1R}+p_2g_{2R})h_{iP}, 


\end{equation} 


а $c^{(2)}_{iQ}$ --- коэффициент при $z_i Z^Q$ в разложении функции 


\begin{equation} 


y_i f_i=z_i(1+h^{(i)})f_i(z_1(1+h^{(1)}), 


z_2(1+h^{(2)})). 


\end{equation} 


Поскольку разложение функции $z_i h^{(i)}$ не содержит свободного и 


линейных слагаемых, имеем 


\begin{equation} 


h_{1(-1,0)}=h_{2(0,-1)}=h_{1(-1,1)}=h_{2(1,-1)}=0. 


\end{equation} 


С помощью последнего соотношения в (5.4) и соотношений (5.5)--(5.7) по 


индукции нетрудно доказать, что 


\begin{equation} 


h_{1(-1,k)}=h_{2(k,-1)}=0\ \ (k\ge0). 


\end{equation} 


С помощью (5.5), (5.8), (5.2), (5.6) и первого из соотношений (5.4) 


теперь удается вычислить $c^{(1)}_{i(k,0)}$ и $c^{(2)}_{i(k,0)}$. Имеем 


$$ 


c^{(1)}_{i(k,0)}=0,\ \ c^{(2)}_{1(k,0)}=2m^{-1}\delta_{km},\ \ 


c^{(2)}_{2(k,0)}=(2m^{-1}-1)\delta_{km}, 


$$ 


$\delta_{km}$ --- символ Кронекера. Из третьего соотношения в (5.4) 


вытекает $g_{1(k,0)}=2m^{-1}\delta_{km}$, 


$g_{2(k,0)}=(2m^{-1}-1)\delta_{km}$. Следовательно, нормальная форма (5.3) 


системы (5.1) имеет вид 


\begin{equation} 


z_{1\tau}=2m^{-1}z^{m+1}_1,\quad 


z_{2\tau}=\lambda_2z_2+(2m^{-1}-1)z^m_1z_2. 


\end{equation} 


Система (5.9) легко интегрируется. Найдем 


\begin{equation} 


z_2=c z^\alpha_1\exp(-1/(4z^m_1)), 


\end{equation} 


$\alpha=\frac{m}{2}(2m^{-1}-1)$.


После подстановки (5.10) в (3.1) получим для однопараметрического 


семейства решений уравнения (2.4) ($c$ --- параметр семейства) следующее 


представление: 


\begin{gather} 


\phi=y_1=z_1\bigg(1+\sum_Q h_{1Q}Z^Q\bigg),\\ 


\sigma=y_2-2m^{-1}y_1=z_2-2m^{-1}z_1+\sum_Q h_{2Q}z_2Z^Q- 


2m^{-1}\sum_Q h_{1Q}z_1Z^Q,\notag 


\end{gather} 


где $z_2$ определяется равенством (5.10). С учетом первого соотношения из 


(5.4) и (5.7) нетрудно заметить, что степенные ряды в правых частях 


равенств (5.11) содержат экспоненциально убывающий при $z_1\to+0$ 


множитель $z_2$ как минимум в первой степени. Выписав лишь слагаемые с 


первой степенью $z_2$, получим асимптотические представления при 


$z_1\to+0$ 


\begin{gather} 


\phi=z_1\bigg(1+\sum^p_{q=0}h_{1(q,1)}z^q_1z_2\bigg)+ 


O(z^{p+2+\alpha}_1\exp(-1/(4z^m_1))),\\ 


\sigma=-2m^{-1}z_1-\bigg(2m^{-1}\sum^p_{q=0} 


h_{1(q,1)}z^{q+1}_1-1\bigg)z_2+ 


O(z^{p+2+\alpha}_1\exp(-1/(4z^m_1))).\notag 


\end{gather} 


При $z_1\to+0$ имеем $\phi\to+0$, $\sigma\to-0$ а $w\to+0$ при $c>0$. 


Кривые (5.11) при $z_1\to+0$ входят в начало координат в секторе $I$ 


фазовой плоскости $(\phi,\sigma)$ в направлении с угловым коэффициентом 


$k_1=-2m^{-1}$. Кривая $S$ получается из семейства кривых (5.11) при 


некотором значении параметра $c=c_1$, зависящем от $m$ и определяемым для 


каждого конкретного значения $m$ лишь численно. Для коэффициентов 


$h_{1(q,1)}$ из (5.4)--(5.6) находим 


$$ 


h_{1(q,1)}=2(c^{(1)}_{1(q,1)}+c^{(2)}_{1(q,1)})= 


2(m^{-1}(5m-2q-2)h_{1(q-m,1)}-\delta_{qm-1}), 


$$ 


отсюда 


\begin{equation} 


h_{1(q,1)}=0\ \ (q<m-1),\ \ h_{1(m-1,1)}=-2,\ \ h_{1(q,1)}= 


2m^{-1}(5m-q-2)h_{1(q-m,1)}\ \ (q\ge m). 


\end{equation} 


Из (5.13) следует, что отличны от нуля лишь коэффициенты 


$$ 


h_{1(m-1,1)}=-2,\quad h_{1(sm-1,1)}=(-1)^{s-1} 


2^s(2s-5)!!\ \ (s\ge2). 


$$ 


Соотношения (5.12) приобретают теперь окончательный вид 


\begin{gather} 


\phi=z_1-2z^m_1z_2-\sum^n_{s=2}(-2)^s(2s-5)!! z^{sm}_1z_2+ 


O(z^{(n+1)m+\alpha}_1\exp(-1/(4z^m_1))),\\ 


\sigma=-2m^{-1}z_1+z_2+4m^{-1}z^m_1z_2+2m^{-1}\sum^n_{s=2} 


(-2)^s(2s-5)!! z^{sm}_1z_2+O(z^{(n+1)m+\alpha}_1 


\exp(-1/(4z^m_1))).\notag 


\end{gather} 


После подстановки в (5.14) $z_2$ из (5.10) с $c=c_1$ получим асимптотику 


кривой $S$ вблизи начала, с помощью которой из второго уравнения (2.3) 


находится сначала асимптотика $\phi(\xi)$, а затем с помощью первого 


равенства (2.3) --- и асимптотика решения краевой задачи (2.1), (2.2) при 


$\xi\to\infty$. Приведем трехчленную асимптотику $\underline{f(\xi)}$. 


Будем иметь


$$


\xi/c=z^{-m/2}_1(1+2^{-1}c_1m\Gamma(-2^{-1},(4z^m_1)^{-1})+ 


O(\Gamma(-3/2,(4z^m_1)^{-1}))),


$$


$z_1\to+0$, $c>0$ произвольно, 


$\Gamma(v,x)$ --- неполная гамма-функция ([7], с.\,954). Используя 


асимптотику $\Gamma(v,x)$ при $x\to\infty$ ([7], с.\,956), найдем 


\begin{equation} 


\xi/c=z^{-m/2}_1(1+4c_1m z^{3m/2}_1\exp(-1/(4z^m_1))+ 


O(z^{5m/2}_1\exp(-1/(4z^m_1)))),\ \ z_1\to+0. 


\end{equation} 


Из (5.15) следует, что $\xi\to\infty$ при $z_1\to+0$. С помощью (5.15) 


находится при $\xi\to\infty$ асимптотика 


\begin{equation} 


z_1=(\xi/c)^{-2/m}\{1+8c_1(\xi/c)^{-3}\exp(-(\xi/2c)^2)+ 


O(\xi^{-5}\exp(-(\xi/2c)^2))\}. 


\end{equation} 


В результате подстановки (5.16) в (5.14) получаем трехчленную асимптотику 


$\phi(\xi)$ при $\xi\to\infty$ 


$$ 


\phi(\xi)=(\xi/c)^{-2/m}\{1-2c_1(c\xi^{-1}-4(c\xi^{-1})^3+ 


O(\xi^{-5}))\exp(-(\xi/2c)^2)\}, 


$$ 


а затем с помощью первого соотношения (2.3) --- асимптотику 


однопараметрического семейства ($c>0$ --- параметр семейства) 


решений уравнения (2.1) 


$$ 


f(\xi,c)=c^{2/m}\{1-2c_1(c\xi^{-1}-4(c\xi^{-1})^3+ 


O(\xi^{-5}))\exp(-(\xi/2c)^2)\}. 


$$ 


С учетом граничного условия (2.2) при $\xi=\infty$ окончательно находим 


\begin{equation} 


f(\xi)=1-2c_1\xi^{-1}(1-4\xi^{-2}+O(\xi^{-4})) 


\exp(-(\xi^2/4)),\quad \xi\to\infty. 


\end{equation} 


Асимптотика производной $f_\xi$ ($\xi\to\infty$) получается 


дифференцированием асимптотики (5.17), и мы ее не приводим. Положив в 


правой части равенства (5.17) $\xi=x t^{-1/2}$, получаем трехчленную 


асимптотику задачи П.Я.\,Полубариновой-Кочиной (1.1), (1.2) при 


$x\to\infty$ ($t>0$ фиксировано). 


 


\begin{notenonum} 


Асимптотика $f(\xi)$ при $\xi\to0$ ($\xi\to\infty$) одновременно дает и 


асимптотику при $x\to0$ ($x\to\infty$), $t>0$ фиксировано, решения 


$u(t,x)$ задачи (1.1), (1.2) и равномерную по $x$ на любом конечном 


интервале полуоси $x\ge0$ асимптотику этого решения при $t\to\infty$ 


($t\to+0$). 


\end{notenonum} 


\addtocounter{section}{1} 


\setcounter{equation}{0}


 


{\bf 6.} В этом пункте для сокращения записи формул будем использовать 


обозначения\linebreak $\alpha=(m+1)/(km-1)$, $\beta=(m+1)^{-1}$, 


$\gamma=(k+1)/(km-1)$, $\delta=(k-1)m$, $\zeta=(km-1)^{-1}$ и 


предполагать, что $\zeta>0$, $k>1$, $m>0$ ($k$, $m$ целые). При 


турбулентной фильтрации политропного газа в пористой среде его плотность 


$\rho(x,t)$ (в случае плоских волн) удовлетворяет уравнению (модель 


Лейбензона) [8] 


\begin{equation} 


\rho_t=(\rho^\delta \rho_x|\rho_x|^{m-1})_x. 


\end{equation} 


Построим возрастающее решение уравнения (6.1), удовлетворяющее при $x>0$ 


условию непрерывности потока $q=\rho^\delta \rho_x|\rho_x|^{m-1}$ и 


условиям Полубариновой-Кочиной (1.2). Это решение автомодельно и имеет вид 


$\rho=f(\xi)$, где $\xi=x t^{-\beta}$, а $f(\xi)$ удовлетворяет уравнению 


\begin{equation} 


(f^{(k-1)m}f^m_\xi)_\xi+\beta\xi f_\xi=0 


\end{equation} 


и граничным условиям (2.2). Положив 


\begin{equation} 


f=\xi^\alpha\phi,\quad \sigma=\xi \phi_\xi, 


\end{equation} 


из (6.2) получим динамическую систему 


\begin{gather} 


\sigma_\tau=Q(\phi,\sigma),\quad \phi_\tau=P(\phi,\sigma),\\ 


P=-m\phi^\delta w^{m-2}\sigma,\ \ Q=m\gamma\phi^\delta w^{m-1}+ 


\delta\phi^{\delta-1}\sigma w^{m-1}+m\alpha\phi^\delta\sigma 


w^{m-2}+\beta,\notag\\ 


w=\sigma+\alpha\phi.\notag 


\end{gather} 


 


В силу очевидного соотношения $f_\xi=\xi^{\alpha-1}w$ и возрастания $f$ 


интегральные кривые уравнения (6.4) следует рассматривать лишь в секторе 


$I$ --- пересечении полуплоскостей $\phi>0$, $w>0$. Снова найдем, что в 


секторе $I$ имеются семейство $S_1$ интегральных кривых с вертикальной 


асимптотой (осью $\sigma$) и семейство $S_2$ с наклонной асимптотой $w=0$. 


Эти семейства разделены кривой $S$ --- сепаpатрисой седла на $\infty$ в 


направлении с угловым коэффициентом $k_1=v_0=-\gamma/k$. Все интегральные 


кривые в секторе $I$ начинаются в точках прямой $w=0$. В окрестности точки 


выхода $(\phi_0,\sigma_0)$, $\phi_0>0$, $\sigma_0=-\alpha\phi_0$, 


сепаратрисы $S$ для нее справедлива асимптотика при $\phi\to\phi_0+0$ 


\begin{gather} 


\sigma-\sigma_0\sim A(\phi-\phi_0)^\eta,\quad 


A=(\beta/(\alpha\eta m\phi_0^{\delta+1}))^\eta,\quad 


\eta=1/(m-1),\ \ m>2,\\ 


\eta=1,\quad A=(\beta-\alpha^2m\phi_0^{\delta+1})/ 


(m\alpha\phi_0^{\delta+1}),\ \ m=2,\notag \\ 


\eta=1,\ \ A=-\alpha,\ \ m<2.\notag 


\end{gather} 


Координаты $\phi_0$, $\sigma_0$ точки выхода могут быть определены лишь 


численно. Однопараметрическое семейство решений $f(\xi,c)$ уравнения 


(6.2), удовлетворяющих граничному условию (2.2) при $\xi=0$ снова 


порождается в силу уравнений (6.3) сеператрисой $S$. С помощью 


преобразования Пуанкаре 


\begin{equation} 


\phi=1/y_1,\quad \sigma=(v_0+y_2)/y_1, 


\end{equation} 


где $v_0=-\gamma/k$, из (6.4) получим динамическую систему 


\begin{equation} 


\begin{split} 


y_{1\tau}&=\lambda_1y_1+m y_1y_2(y_2+k^{-1})^{m-2}+m v_0y_1 


\{(y_2+k^{-1})^{m-2}-k^{2-m}\},\\ 


y_{2\tau}&=\lambda_2y_2+m k(y_2)^2(y_2+k^{-1})^{m-2}+m y_2 


\{(y_2+k^{-1})^{m-2}-k^{2-m}\}+\beta(y_1)^{1/\zeta}, 


\end{split} 


\end{equation} 


где $\lambda_1=m v_0k^{2-m}<0$, $\lambda_2=m k^{2-m}>0$. Преобразование 


(6.6) переводит седло на $\infty$ динамической системы (6.4) в седло 


$y_1=y_2=0$ системы (6.7), а кривую $S$ --- в одну из ветвей сепаратрисы 


этого седла. Рассуждая, как и в п.\,3, получим для искомой ветви 


сепаратрисы седла $y_1=y_2=0$ динамической системы (6.7) представление 


\begin{equation} 


y_1=u^\zeta\bigg(1+\sum^\infty_{s=1}a_s u^s\bigg),\quad 


y_2=\sum^\infty_{s=1}b_s u^s,\quad u=\exp(\lambda_1\tau/\zeta). 


\end{equation} 


Коэффициенты $a_s$, $b_s$ определяются последовательно ($a_s$ определяется 


вслед за $b_s$) с помощью рекуррентных соотношений, которые получаются в 


результате подстановки рядов (6.8) в уравнение (6.7) с учетом соотношений 


$$ 


y_{1\tau}=\lambda_1y_1+\lambda_1\zeta^{-1}u^\zeta 


\sum^\infty_{s=1}s a_s u^s,\quad 


y_{2\tau}=\lambda_1\zeta^{-1}\sum^\infty_{s=1}s b_s u^s 


$$ 


и сравнения коэффициентов при одинаковых степенях $u$ в полученных 


равенствах. Ввиду громоздкости мы не приводим этих рекуррентных 


соотношений и ограничимся вычислением коэффициентов 


\begin{equation} 


a_1=m\zeta\lambda_1^{-1}k^{3-m}b_1(k^{-1}+(m-2)v_0),\quad 


b_1=\beta(\lambda_1\zeta^{-1}-\lambda_2)^{-1}. 


\end{equation} 


С помощью представления (6.8) для $y_1$, $y_2$ и равенств (6.6) из (6.3) 


получаем для однопараметрического семейства $f(\xi,c)$ решений уравнения 


(6.2), удовлетворяющих граничному условию (2.2) при $\xi=0$, представление 


в виде ряда по дробным степеням автомодельной переменной $\xi$ 


\begin{equation} 


f(\xi,c)=c^\alpha(\xi/c)^{1/k}\sum^\infty_{s=0} 


l_s(\xi/c)^{s(k+1)/k}. 


\end{equation} 


Отметим, что показатели степени в разложении (6.10) не зависят от $m$, а


коэффициенты $l_s$ от $m$ зависят и выражаются через коэффициенты рядов 


(6.8). Эти выражения весьма громоздки, и мы ограничимся тем, что приведем 


лишь коэффициенты 


$$ 


l_0=1,\quad l_1=k b_1/(k+1), 


$$ 


$a_1$, $b_1$ здесь определяются равенствами (6.9). Значение параметра $c$ 


в равенстве (6.10) определяется вторым граничным условием (2.2). С учетом 


асимптотики (6.5) нетрудно убедиться с помощью соотношений (6.3), что для 


каждого решения семейства (6.10) существует точка $\xi_1(c)$ такая, что 


$f_\xi(\xi_1(c),c)=0$, причем $f_\xi(\xi,c)>0$ при $\xi<\xi_1(c)$. 


Нетрудно также убедиться, что вместе с решением $f(\xi)$ уравнения (6.2) 


его решением является и функция $\mu^{-\alpha}f(\mu\xi)$, где $\mu>0$ 


произвольно. Рассмотрим решение $f(\xi,c)$ при $c=1$ и обозначим его 


$f_1(\xi)$, а через $\xi_1=\xi_1(1)$ обозначим точку, в которой обращается 


в нуль $f_{1\xi}$ (точка $\xi_1$ при заданных $k$ и $m$ определяется 


численно). Положим $f_0=f_1(\xi_1)>0$. Искомое решение $f(\xi,c_0)$ 


уравнения (6.2), удовлетворяющее обоим граничным условиям (2.2), будем 


искать в виде $f(\xi,c_0)=\mu^{-\alpha}f_1(\mu\xi)$, его производная 


обращается в нуль в точке $\xi_1(c_0)$. Определим значение $\mu$ из 


условия $f(\xi_1(c_0),c_0)=1$. Будем иметь $\xi_1(c_0)=\xi_1/\mu$, причем 


$\mu=f_0^{1/\alpha}$. Таким образом, решение 


$f=f_0^{-1}f_1(f_0^{1/\alpha}\xi)$ уравнения (6.2) удовлетворяет 


граничному условию (2.2) при $\xi=0$, монотонно возрастает на интервале 


$(0,\xi_1f_0^{-1/\alpha})$, причем $f(\xi_1f_0^{-1/\alpha})=1$. Положив 


$f(\xi)=1$ при $\xi>\xi_1f_0^{-1/\alpha}$, получим решение уравнения 


(6.2), удовлетворяющее обоим граничным условиям (2.2) и непрерывное вместе 


с производной $f_\xi$, следовательно, удовлетворяющее и условию 


непрерывности потока $q$. С учетом равенства (6.10) для этого решения 


получим окончательно представление 


\begin{gather*} 


f=f_0^{-1}(\xi f_0^{1/\alpha})^{1/k}\sum^\infty_{s=0}l_s 


(\xi f_0^{1/\alpha})^{s(k+1)/k}\ \ \text{при} 


\ \ 0\le\xi\le\xi_1f_0^{-1/\alpha},\\ 


f=1\ \ \text{при}\ \ \xi>\xi_1f_0^{-1/\alpha} 


\end{gather*} 


(значение параметра $c$, выделяющее искомое решение из семейства (6.10), 


есть $c_0=f_0^{-1/\alpha}$). В заключение отметим, что нетрудно вычислить 


и асимптотику $f(\xi)$ при $\xi\to\xi_0-0$, $\xi_0=\xi_1f_0^{-1/\alpha}$. 


С помощью (6.5) и (6.3) без труда получаем двучленную асимптотику $f(\xi)$ 


при $\xi\to\xi_0-0$ в виде 


$$ 


f(\xi)=1-B(\xi_0-\xi)^{1+\eta}+o((\xi_0-\xi)^{1+\eta}), 


$$ 


где $B=\alpha^\eta(1+\eta)^{-1}A \xi_0^{a(1-\eta)-1-\eta}$ при $m>2$ и 


$\eta=1$, $B=(12\xi_0)^{-1}$ при $m=2$ (отметим, что $\xi_0$ и $\phi_0$ 


связаны простым соотношением $\xi^\alpha_0\phi_0=1$, являющимся следствием 


условия $f(\xi_0)=1$). 


 


\begin{thebibliography}{9} 


 


\bibitem{1} 


Полубаринова-Кочина П.Я. {\em Об одном нелинейном уравнении в частных 


производных, встречающемся в теории фильтрации} // ДАН СССР. -- 1948. -- 


Т.\,63. -- \No\,6. -- С.\,623--626. 


\bibitem{2} 


Качан М.В., Пименов С.Ф., Сущий С.М., Энтель М.Б. {\em О некоторых 


автомодельных решениях уравнения Лейбензона} // Механ. жидкости и газа. 


-- 1991. -- \No\,5. -- С.\,145--150. 


\bibitem{3} 


Баренблатт Г.И. {\em О некоторых неустановившихся движениях жидкости и 


газа в пористой среде} // ПММ. -- 1952. -- Т.\,16. -- Вып.\,1. 


-- С.\,67--78. 


\bibitem{4} 


Самарский А.А., Галактионов В.А., Курдюмов С.П., Михайлов А.П. {\em 


Режимы с обострением в задачах для квазилинейных параболических 


уравнений}. -- М.: Наука, 1987. -- 477~с. 


\bibitem{5} 


Баутин Н.Н., Леонтович Е.А. {\em Методы и приемы качественного 


исследования динамических систем на плоскости}. -- М.: 


Наука, 1990. -- 488~с. 


\bibitem{6} 


Брюно А.Д. {\em Локальный метод нелинейного анализа дифференциальных 


уравнений}. -- М.: Наука, 1979. -- 253~с. 


\bibitem{7} 


Градштейн И.С., Рыжик И.М. {\em Таблицы интегралов, сумм, рядов и 


произведений}. -- М.:\linebreak ГИФМЛ, 1962. -- 1100~с. 


\bibitem{8} 


Баренблатт Г.И. {\em Об автомодельных движениях сжимаемой жидкости в 


пористой среде} // ПММ. -- 1952. -- Т.\,16. -- Вып.\,6. -- С.\,679--698. 


 


\end{thebibliography} 


 


\vskip 2cm 


 


\post{\it Казанский государственный}{\it Поступила} 


\post{\it университет}{24.05.2002} 


 


\label{end} 


\end{document} 


